\section{Introducción y marco teórico}
Este trabajo aspira a simular el Sistema Solar con datos realistas, para lograr recrear una 
maniobra de asistencia gravtitatoria, 
con datos y trayectoria similares a los de la nave \textbf{Solar Orbiter} de la Agencia Espacial Europea (ESA).

Una asistencia gravitatoria consiste en el intercambio de energía, por medio de la variación del momento angular, 
entre un cuerpo celeste y una nave.
El objetivo del presente trabajo es intentar maximizar dicho intercambio de energía (que tendrá que ser negativo) 
para acercarse lo 
máximo posible al Sol.

\subsection{Parámetros iniciales}
Los parámetros físicos fundamentales utilizados en la simulación han sido extraídos de fuentes oficiales para 
garantizar la precisión de los resultados. De la base de datos del \textit{Jet Propulsion Laboratory} (JPL) de 
la NASA\footnote{\texttt{https://ssd.jpl.nasa.gov/astro\_par.html}}, han sido extraídos el valor de la 
constante de gravitación universal $G$ y los 
parámetros gravitacionales estandarizados de los distintos cuerpos. Dado que estos sistemas proporcionan 
las masas en la forma del parámetro gravitacional $G M$, las masas individuales,
$M_\circ$, $M_T$, $M_J$, $M_m$, $M_v$, han sido obtenidas dividiendo los valores por $G$. 
Para ciertos parámetros orbitales específicos 
y radios, como los de la Luna, se han consultado fuentes como la \textit{Planetary Fact Sheet} de la 
NASA\footnote{\texttt{http://www.mathaddict.net/ellipseplanet.pdf}} y 
tablas de mecánica celeste estandarizadas\footnote{\texttt{https://www.mnealon.eosc.edu/NewtonianMechanics.htm}}.

Así, utilizando los datos aportados por el sistema \textit{Horizons} del JPL de la NASA, es posible simular 
los planetas en su posición real el día en el que se lanzó el \textit{Solar Orbiter}. Dado que las coordenadas 
de la nave no están disponibles en el momento exacto de lanzamiento (pues está en la Tierra), se toman las 
coordenadas de los planetas a las 00:00 del 11 de febrero de 2020, unas 20 horas después del lanzamiento.

\subsection{Condiciones iniciales de la nave}
Para obtener unas condiciones iniciales realistas para la nave, se recurre de nuevo al sistema 
\textit{Horizons} del JPL de la NASA para extraer tanto el vector de posición como el vector de velocidad 
del \textit{Solar Orbiter} en ese instante, a las 00:00 del 11 de febrero de 2020. 

% Al transcurrir este tiempo desde el lanzamiento, la sonda ya ha establecido su trayectoria de 
% transferencia interplanetaria inicial, alejándose de la Tierra para escapar de su Esfera de Influencia, 
% situándose a una distancia que aproximaremos a $\frac{49}{50}$ UA respecto a la posición original, 
% o unos 3 millones de kilómetros, donde la perturbación gravitatoria terrestre es ya residual y 
% no domina la dinámica orbital en favor del Sol.

Esto permite fijar un módulo de velocidad inicial y una posición de partida que son físicamente realistas 
y coherentes con la capacidad de los lanzadores de la misión real; sin embargo, si se utilizara 
también la dirección exacta del vector velocidad proporcionado por el JPL, el trabajo quedaría reducido a una mera 
recreación de la trayectoria histórica. 
Debido a ello, dicha dirección será obviada y se planteará un problema de optimización:
encontrar la dirección de lanzamiento óptima que minimice la distancia de aproximación a Venus para así maximizar
la transferencia de energía entre el planeta y la nave.

\subsubsection{Justificación Analítica de la Asistencia Gravitatoria}
La relación entre la distancia de aproximación (periapsis, $r_p$) y el intercambio de energía se puede justificar analíticamente. 
En el sistema de referencia planetario, la nave describe una trayectoria hiperbólica, 
que no es más que la solución matemática de un problema de dispersión gravitatoria.
Es análogo a la dispersión de Rutherford, siendo la única diferencia que tratamos con un potencial atrctivo en vez de repulsivo
(potencial $V(r) \propto +1/r$ en vez de $V(r) \propto -1/r$).
Esto resulta en un cambio de signo en la ecuación de la trayectoria, siendo la de la atractiva:
$$r = \frac{p}{1 + \varepsilon\cos \beta}$$
Sin embargo, la relación final entre el ángulo de 
dispersión $\theta$ y la excentricidad se mantiene idéntica en ambos casos:
$$\sin\left(\frac{\theta}{2}\right) = \frac{1}{\varepsilon}$$

\paragraph{Excentricidad en función de la periapsis.}
La excentricidad de una órbita en un campo de fuerzas centrales con potencial $V = -k/r$ viene dada por la expresión general:
\begin{equation}
    \varepsilon = \sqrt{1 + \frac{2 E J^2}{m k^2}}
\end{equation}
donde $E$ es la energía mecánica total, $J$ el momento angular total y $m$ la masa de la nave. Para el potencial gravitatorio, el factor $k$ es $k = G M_V m$, donde $M_V$ es la masa de Venus.
\newpage
Sustituyendo $k$ y las expresiones para la energía lejos del planeta ($E = \frac{1}{2}mv_\infty^2$) y el momento angular al cuadrado ($J^2 = m^2(r_p^2 v_\infty^2 + 2 G M_V r_p)$), obtenemos:
\begin{equation}
    \varepsilon^2 = 1 + \frac{2 \left(\frac{1}{2}mv_\infty^2\right) \left(m^2(r_p^2 v_\infty^2 + 2 G M_V r_p)\right)}{m(m G M_V)^2}
\end{equation}
Simplificando los términos de masa $m$:
\begin{align}
    \varepsilon^2 &= 1 + \frac{m^3 v_\infty^2 (r_p^2 v_\infty^2 + 2 G M_V r_p)}{m^3 (G M_V)^2} \nonumber \\
    &= 1 + \frac{v_\infty^2 (r_p^2 v_\infty^2 + 2 G M_V r_p)}{(G M_V)^2} \nonumber \\
    &= 1 + \frac{r_p^2 v_\infty^4}{(G M_V)^2} + \frac{2 r_p v_\infty^2}{G M_V} = \left(1 + \frac{r_p v_\infty^2}{G M_V}\right)^2
\end{align}
De donde se obtiene la relación fundamental, que es independiente de la masa de la nave ($m$):
\begin{equation}
    \varepsilon = 1 + \frac{r_p v_\infty^2}{G M_V}
\end{equation}

\paragraph{Ángulo de giro y cambio de velocidad.}
Al sustituir la expresión de $\epsilon$ en la ecuación 1, se obtiene la primera fórmula clave:
\begin{equation}
    \sin\left(\frac{\theta}{2}\right) = \frac{1}{1 + \dfrac{r_p v_\infty^2}{G M_V}}
\end{equation}
El cambio de velocidad, $\Delta v$, se obtiene por geometría vectorial. 
Dado que la energía se conserva, el módulo de la velocidad de entrada y de salida es el mismo ($v_\infty$). 
El vector $\Delta \vec{v}$ forma la base de un triángulo isósceles, cuyo módulo es 
$\Delta v = 2 v_\infty \sin(\theta/2)$. Combinando ambas ecuaciones, se obtiene la expresión final:
\begin{equation}
    \Delta v = \frac{2 v_\infty}{1 + \dfrac{r_p v_\infty^2}{G M_V}}
\end{equation}
Como se observa, para maximizar el cambioo de velocidad $\Delta v$ y maximizar la transferencia 
energética, se debe minimizar la distancia de periapsis $r_p$.
Sin embargo, esta optimización cuenta con una restricción física indispensable: 
para que la sonda pueda continuar su misión, no debe colisionar con la superficie de Venus 
ni desintegrarse por la fricción aerodinámica con la atmósfera (ni perder energía por esta). 
\newpage
Por tanto, la distancia mínima de aproximación debe cumplir siempre la condición límite de que 
$r_p > R_V + h_{\text{atm}}$, siendo $R_V$ el radio del planeta y $h_{\text{atm}}$ el 
límite superior de su atmósfera.
Añadiremos, además, una holgura extra para evitar cualquier tipo de pérdida energética por un posible
rozamiento con la atmósfera; el radio de Venus es de 6052km según las fuentes  
nombradas anteriormente, y su exosfera se encuentra a 350km.

Además, aunque el Sistema Solar interior es fundamentalmente coplanario, las órbitas presentan ligeras inclinaciones respecto a otras; 
en nuestro caso concreto, la órbita de Venus tiene una inclinación de aproximadamente $3.4^\circ$ respecto a la eclíptica. 
Si cada simulación estuviera limitara a un único grado de libertad variando solo el ángulo de guiñada, $\varphi$,
la nave quedaría contenida en el mismo plano para todas las simulaciones.

Debido a esto, se realizará un barrido sobre dos grados de libertad: el ángulo de guiñada (yaw), $\varphi$, y el de cabeceo
(pitch), $\theta$.
El hecho de mantener fijo el módulo de la velocidad a uno realista, permite acotar el coste computacional de las simulaciones,
y restringir el mapa de soluciones a uno bidimensional, facilitando así su interpretación.